\begin{table}[htbp]
  \centering
  \caption{Pontos críticos identificados no projeto de exemplo, agrupados pela função que os contém}
  \label{tab:pontos}
  \begin{tabular}{p{3.0cm}p{6.6cm}cc}
    \toprule
    \textbf{Arquivo} & \textbf{Função e tipos de ponto crítico} & \textbf{Pontos} & \textbf{Gravidade} \\
    \midrule
    \texttt{persistencia.py} & salvar(self, funcionario, setor) & 2 & 5 \\
     & \footnotesize dependência externa, tratamento de exceção & & \\
    \texttt{relatorios.py} & contratar\_em\_lote(pedidos, repositorio) & 1 & 4 \\
     & \footnotesize tratamento de exceção & & \\
    \texttt{regras\_negocio.py} & validar\_contratacao(funcionario, setor, empresa) & 2 & 3 \\
     & \footnotesize condicional aninhado, literal em condicional & & \\
    \texttt{relatorios.py} & setores\_sobrecarregados(empresas) & 1 & 3 \\
     & \footnotesize laço aninhado & & \\
    \bottomrule
  \end{tabular}
  \fonte{Elaborado pelo autor, a partir da saída do protótipo.}
\end{table}
